\section{模型评价与推广}

\subsection{模型评价}

本文模型的主要特点不是为四个问题分别建立四套互不相关的算法，而是在同一条物理主线上逐步增加现实因素。母线能量平衡、储能状态递推、容量和功率边界始终保持不变；问题一解决确定信息下“电量在时间上如何搬移”，问题二增加预测误差与紧急购电风险，问题三允许新预报到达后修订尚未执行的计划，问题四再把固定电价替换为逐时波动价格。这样的结构使各问之间既有继承关系，又能够明确判断费用变化究竟来自物理条件、信息条件还是结算规则。

在不确定性处理上，本文没有把当天真实数据提前放入计划层，而是显式区分“决策时已知信息”和“执行后实现信息”。问题二用历史数据形成日前预测，并利用 $5$ 倍紧急电价推导出的非对称损失设计光伏风险边际；问题三把每次新预报与当时实测储电量纳入下一滚动阶段。配合反事实扰动检验，这一处理使较低的购电费用对应于真正可执行的策略，而不是未来信息带来的理想化结果。

模型还具有较好的可解释性。问题一能够直接从“低价充电、高价放电”和“光伏先存后用”解释储能收益；问题二通过信息基准、风险边际和实时反馈的对照分别说明三个机制的作用；问题三用逐节点消融回答“哪些预报值得使用”，并进一步把综合节点收益拆分为已实现负载刷新和新光伏预报的顺序增量；问题四则把题面给定价格下的正式结果与因果价格扩展分开。这样，最终结论不仅给出费用数字，也说明数字为什么发生变化。

模型的局限同样需要明确。问题二和问题三采用日循环终端策略，避免了单日优化在午夜透支储能，但同时削弱了跨日套利空间；历史预测器没有使用辐照度、云量、温度、节假日等外生特征，对突发天气和异常负荷的响应仍有限；当前光伏风险边际按预测器使用统一的残差分位，尚未进一步描述不同季节、不同日内时刻的异方差；储能模型也没有计入自放电、容量衰减和循环寿命成本。因此，本文更适合作为“已知储能参数下的短期经济调度模型”，其全年费用不应直接解释为计入全部设备寿命成本后的长期财务最优值。

\subsection{模型改进}
\label{sec:improve}

后续改进可围绕“更长时间尺度、更完整的不确定性、更真实的设备模型”展开。首先，可把单日优化扩展为 $3\sim7$ 天滚动时域，只执行首日方案，并在窗口末端加入储能残值，从而允许跨日转移电量而不产生末端透支。其次，可利用数值天气预报、辐照度、温度和工作日特征建立负荷--光伏联合概率场景，再采用情景抽样或 CVaR 直接优化跨时段风险，使单时段 $0.2$ 分位启发上升为完整的随机调度模型。再次，可按季节和日内时刻分别估计预测残差分位数，使风险边际随预测不确定性动态变化。最后，可将电池循环寿命、自放电、功率相关效率和维护成本纳入目标函数，使经济调度与设备长期健康状态统一起来。

\subsection{模型推广}

本文“公共物理约束--信息分层--滚动修正”的思路并不依赖于本题的具体数据。在园区综合能源系统中，可以把单一电力平衡扩展为电、热、气等多能流平衡，并增加储热、储气状态 \cite{liu2018dayahead}；在电动汽车充电站中，可将小区负载替换为车辆到站后的充电需求，并加入离站时间和需求满足约束；在工业用户侧储能中，则可进一步考虑最大需量电费与生产计划。更一般地，只要系统具有“先做计划、后观察偏差、状态逐期传递、允许后续修正”的结构，例如库存补货、水库调度和供应链运输，都可以复用本文的信息边界与滚动决策框架。
